\section{Class NP --- Poly Time Block Loop}
%%% generated by the blueprint pipeline from TCSlib.Complexity.ClassNP.PolyTimeBlockLoop
\subsection*{Class NP --- Poly Time Block Loop: Introduction}

Declarations of the Lean module \texttt{TCSlib.Complexity.ClassNP.PolyTimeBlockLoop}.

\subsection*{Class NP --- Poly Time Block Loop: Results}

\begin{definition}[Encoded prefix slice]
\lean{Complexity.sliceTakeAt}
\label{Complexity.sliceTakeAt}
\leanok
\uses{Complexity.pairFstD, Complexity.pairSndD, Turing.pairEncode}
Given natural numbers $a,k$ and an encoded pair of Boolean lists $z$, decode its first
and second components as $x$ and $y$. Take the first $a(|x|+1)^k$ entries of $y$, then
encode the pair consisting of $x$ and that prefix.
\end{definition}

\begin{definition}[Encoded suffix slice]
\lean{Complexity.sliceDropAt}
\label{Complexity.sliceDropAt}
\leanok
\uses{Complexity.pairFstD, Complexity.pairSndD, Turing.pairEncode}
Given natural numbers $a,k$ and an encoded pair of Boolean lists $z$, decode its first
and second components as $x$ and $y$. Remove the first $a(|x|+1)^k$ entries of $y$, then
encode the pair consisting of $x$ and the remaining suffix.
\end{definition}

\begin{theorem}[Polynomial-time computation of a paired prefix slice]
\lean{Complexity.polyTimeComputable_sliceTakeAt}
\label{Complexity.polyTimeComputable_sliceTakeAt}
\leanok
\uses{Complexity.PolyTimeComputable, Complexity.PolyTimeComputable.comp, Complexity.PolyTimeComputable.pairEncode, Complexity.PrefixByLength.take, Complexity.pairFstD, Complexity.pairFstD_pairEncode, Complexity.pairSndD, Complexity.pairSndD_pairEncode, Complexity.polyTimeComputable_pairFstD, Complexity.polyTimeComputable_pairSndD, Complexity.polyTimeComputable_polyUnary, Complexity.polyTimeComputable_takePrefixByLength, Complexity.sliceTakeAt, Turing.pairEncode}
\difficulty{4}
For every $a,k\in\mathbb{N}$, define a function on bit strings $z$ as follows: parse the
total first and second projections of $z$, using the empty word for either projection
when parsing fails; take the first $a(|x|+1)^k$ bits of the second projection, where $x$
is the first projection; then encode the pair consisting of $x$ and this prefix. This
function is polynomial-time computable.
\end{theorem}

\begin{theorem}[Polynomial-time computation of a paired suffix slice]
\lean{Complexity.polyTimeComputable_sliceDropAt}
\label{Complexity.polyTimeComputable_sliceDropAt}
\leanok
\uses{Complexity.PolyTimeComputable, Complexity.PolyTimeComputable.comp, Complexity.PolyTimeComputable.pairEncode, Complexity.PrefixByLength.drop, Complexity.pairFstD, Complexity.pairFstD_pairEncode, Complexity.pairSndD, Complexity.pairSndD_pairEncode, Complexity.polyTimeComputable_dropPrefixByLength, Complexity.polyTimeComputable_pairFstD, Complexity.polyTimeComputable_pairSndD, Complexity.polyTimeComputable_polyUnary, Complexity.sliceDropAt, Turing.pairEncode}
\difficulty{4}
For every $a,k\in\mathbb{N}$, define a function on bit strings $z$ as follows: parse the
total first and second projections of $z$, using the empty word for either projection
when parsing fails; drop the first $a(|x|+1)^k$ bits of the second projection, where $x$
is the first projection; then encode the pair consisting of $x$ and the remaining
suffix. This function is polynomial-time computable.
\end{theorem}

\begin{theorem}[Polynomial-time tail function]
\lean{Complexity.polyTimeComputable_tail}
\label{Complexity.polyTimeComputable_tail}
\leanok
\uses{Complexity.PolyTimeComputable, Complexity.PolyTimeComputable.comp, Complexity.PolyTimeComputable.pairEncode, Complexity.PrefixByLength.drop, Complexity.polyTimeComputable_const, Complexity.polyTimeComputable_dropPrefixByLength, Complexity.polyTimeComputable_id, Turing.pairEncode}
\difficulty{3}
The function on bit strings that deletes the first bit, returning the empty word on an
empty input, is polynomial-time computable.
\end{theorem}

\begin{theorem}[Testing whether a bit string is empty]
\lean{Complexity.polyTimeComputable_isNil}
\label{Complexity.polyTimeComputable_isNil}
\leanok
\uses{Complexity.PolyTimeComputable, Complexity.PolyTimeComputable.comp, Complexity.PolyTimeComputable.pairEncode, Complexity.pairFstD_pairEncode, Complexity.pairSndD_pairEncode, Complexity.polyTimeComputable_const, Complexity.polyTimeComputable_id, Complexity.polyTimeComputable_lenLe, Turing.pairEncode}
\difficulty{3}
The function sending a bit string $w$ to the one-bit word $1$ if $w$ is empty and to the
one-bit word $0$ otherwise is polynomial-time computable.
\end{theorem}

\begin{theorem}[Polynomial-time closure under Boolean disjunction]
\lean{Complexity.polyTimeComputable_or}
\label{Complexity.polyTimeComputable_or}
\leanok
\uses{Complexity.PolyTimeComputable, Complexity.polyTimeComputable_const, Complexity.polyTimeComputable_ite}
\difficulty{4}
Let $p$ and $q$ be Boolean-valued functions on bit strings. If the one-bit output
functions $x\mapsto p(x)$ and $x\mapsto q(x)$ are both polynomial-time computable, then
the function sending $x$ to the one-bit word $1$ when at least one of $p(x)$ and $q(x)$
equals $1$, and to $0$ otherwise, is polynomial-time computable.
\end{theorem}

\begin{theorem}[Polynomial-time closure under Boolean negation]
\lean{Complexity.polyTimeComputable_not}
\label{Complexity.polyTimeComputable_not}
\leanok
\uses{Complexity.PolyTimeComputable, Complexity.polyTimeComputable_const, Complexity.polyTimeComputable_ite}
\difficulty{4}
Let $p$ be a Boolean-valued function on bit strings, and suppose the function sending
$x$ to the one-bit word $p(x)$ is polynomial-time computable. Then the function sending
$x$ to the one-bit word $1-p(x)$ is polynomial-time computable.
\end{theorem}

\begin{theorem}[Empty-word first projection]
\lean{Complexity.pairFstD_nil}
\label{Complexity.pairFstD_nil}
\leanok
\uses{Complexity.pairFstD}
\difficulty{1}
The total first projection of the empty bit string is the empty bit string.
\end{theorem}

\begin{theorem}[Second projection length bound]
\lean{Complexity.length_pairSndD_le}
\label{Complexity.length_pairSndD_le}
\leanok
\uses{Complexity.pairSndD, Complexity.pairSndD_pairEncode, Turing.eq_pairEncode_of_pairDecode, Turing.length_pairEncode, Turing.pairDecode}
\difficulty{4}
For every finite bit string $z$, the total second projection of $z$ has length at most
$|z|$.
\end{theorem}

\begin{theorem}[Reconstruction from a nonempty first projection]
\lean{Complexity.eq_pairEncode_of_pairFstD_ne}
\label{Complexity.eq_pairEncode_of_pairFstD_ne}
\leanok
\uses{Complexity.pairFstD, Complexity.pairFstD_pairEncode, Complexity.pairSndD, Complexity.pairSndD_pairEncode, Turing.eq_pairEncode_of_pairDecode, Turing.pairDecode, Turing.pairEncode}
\difficulty{4}
Let $z$ be a finite Boolean word whose total first projection under the pairing parser
is nonempty. Then $z$ is the pairing of its total first and second projections, using
the encoding that doubles every bit of the first word, appends the separator $01$, and
then appends the second word.
\end{theorem}

\begin{theorem}[Length bound for decoded pair components]
\lean{Complexity.length_pair_components_le}
\label{Complexity.length_pair_components_le}
\leanok
\uses{Complexity.pairFstD, Complexity.pairFstD_pairEncode, Complexity.pairSndD, Complexity.pairSndD_pairEncode, Turing.eq_pairEncode_of_pairDecode, Turing.length_pairEncode, Turing.pairDecode}
\difficulty{4}
For every bit string $y$, the total first and second projections satisfy
$2\,|\pi_1(y)|+|\pi_2(y)|\le |y|$, where each projection returns the empty word when
parsing fails.
\end{theorem}

\begin{theorem}[Length bound for encoded dropped projections]
\lean{Complexity.length_pairEncode_drop_le}
\label{Complexity.length_pairEncode_drop_le}
\leanok
\uses{Complexity.pairFstD, Complexity.pairFstD_pairEncode, Complexity.pairSndD, Complexity.pairSndD_pairEncode, Turing.eq_pairEncode_of_pairDecode, Turing.length_pairEncode, Turing.pairDecode, Turing.pairEncode}
\difficulty{4}
For every finite bit string $s$ and natural numbers $i,j$, take the total first and
second projections of $s$, respectively, delete the first $i$ and $j$ bits from them,
and encode the resulting pair. The encoded string has length at most $\max(|s|,2)$.
\end{theorem}

\begin{theorem}[Length bound for a dropped slice]
\lean{Complexity.length_sliceDropAt_le}
\label{Complexity.length_sliceDropAt_le}
\leanok
\uses{Complexity.pairFstD, Complexity.pairFstD_pairEncode, Complexity.pairSndD, Complexity.pairSndD_pairEncode, Complexity.sliceDropAt, Turing.eq_pairEncode_of_pairDecode, Turing.length_pairEncode, Turing.pairDecode}
\difficulty{4}
For natural numbers $a$ and $k$ and any bit string $z$, encode the first component of
the total pair decoding of $z$ together with the suffix of its second component obtained
by dropping $a(|\pi_1(z)|+1)^k$ bits. The resulting bit string has length at most
$\max(|z|,2)$.
\end{theorem}

\begin{theorem}[Singleton contribution to a range flat map]
\lean{Complexity.flatMap_range_eq_single}
\label{Complexity.flatMap_range_eq_single}
\leanok
\difficulty{5}
Let $N,K$ be natural numbers, let $b$ be a Boolean list, and let $c$ assign a Boolean
list to each natural number. If $K<N$ and, for every $i<N$, $c(i)=b$ when $i=K$ and
$c(i)$ is empty otherwise, then concatenating the lists $c(i)$ for $i=0,\ldots,N-1$
gives $b$.
\end{theorem}

\begin{definition}[Completed block marker]
\lean{Complexity.blockDone}
\label{Complexity.blockDone}
\leanok
\uses{Turing.pairEncode}
The completed-block marker is the encoding of the pair of empty Boolean lists.
\end{definition}

\begin{definition}[Boolean list emptiness test]
\lean{Complexity.isNilB}
\label{Complexity.isNilB}
\leanok
For a finite Boolean list $w$, return true exactly when $w$ is empty, and false
otherwise.
\end{definition}

\begin{theorem}[Polynomial-time first-bit function]
\lean{Complexity.polyTimeComputable_take1}
\label{Complexity.polyTimeComputable_take1}
\leanok
\uses{Complexity.PolyTimeComputable, Complexity.PolyTimeComputable.comp, Complexity.PolyTimeComputable.pairEncode, Complexity.PrefixByLength.take, Complexity.polyTimeComputable_const, Complexity.polyTimeComputable_id, Complexity.polyTimeComputable_takePrefixByLength, Turing.pairEncode}
\difficulty{4}
The function on bit strings that returns the first bit as a word of length one, or the
empty word when the input is empty, is polynomial-time computable.
\end{theorem}

\begin{theorem}[First bit with a default is polynomial-time]
\lean{Complexity.polyTimeComputable_headD}
\label{Complexity.polyTimeComputable_headD}
\leanok
\uses{Complexity.PolyTimeComputable, Complexity.PolyTimeComputable.append, Complexity.PolyTimeComputable.comp, Complexity.polyTimeComputable_const, Complexity.polyTimeComputable_id, Complexity.polyTimeComputable_take1}
\difficulty{4}
The function sending a bit string $w$ to the one-bit word whose bit is the first bit of
$w$, or $0$ when $w$ is empty, is polynomial-time computable.
\end{theorem}

\begin{theorem}[Polynomial-time concatenation of iterated outputs]
\lean{Complexity.polyTimeComputable_emitIter}
\label{Complexity.polyTimeComputable_emitIter}
\leanok
\uses{Complexity.PolyTimeComputable, Turing.FinTM.exists_emitIterTM}
\difficulty{1}
Let $g,e:
\{0,1\}^*\to\{0,1\}^*$ be polynomial-time computable. Suppose $a',k',b,l\in\mathbb{N}$
and that, for every bit string $w$ and every $i\in\mathbb{N}$, the $i$-fold iterate of
$g$ satisfies $|(g^{i}(w))|\le b(|w|+1)^l$. Then the function sending $w$ to the
concatenation of $e(g^{i}(w))$ over all integers $i$ from $0$ through $a'(|w|+1)^{k'}$
is polynomial-time computable.
\end{theorem}

\begin{definition}[Block loop state transition]
\lean{Complexity.BlockLoop.step}
\label{Complexity.BlockLoop.step}
\leanok
\uses{Complexity.blockDone, Complexity.isNilB, Complexity.pairFstD, Complexity.pairSndD, Turing.pairEncode}
Given update and next-state functions on pairs of Boolean lists and an encoded state
$s$, decode its outer pair. If the first component is empty, return the completed-block
marker. Otherwise, decode the second component; if its first component is empty, also
return the completed-block marker. In every other case, encode a new outer pair: its
first component is the update function applied to the first component of $s$ and the
second component of the decoded inner pair; its second component encodes the inner first
component with its first entry removed, paired with the next-state function applied to
those same inputs.
\end{definition}

\begin{definition}[Block loop output]
\lean{Complexity.BlockLoop.emit}
\label{Complexity.BlockLoop.emit}
\leanok
\uses{Complexity.isNilB, Complexity.pairFstD, Complexity.pairSndD}
Given an output function on finite Boolean lists and an encoded state $s$, return the
empty list if the first component of the decoded outer pair is empty. Otherwise, return
the output function applied to $s$ if the first component of the decoded inner pair is
empty; in all other cases return the empty list.
\end{definition}

\begin{theorem}[Polynomial time for the empty-head test]
\lean{Complexity.BlockLoop.polyTimeComputable_headNil}
\label{Complexity.BlockLoop.polyTimeComputable_headNil}
\leanok
\uses{Complexity.PolyTimeComputable, Complexity.PolyTimeComputable.comp, Complexity.isNilB, Complexity.pairFstD, Complexity.polyTimeComputable_isNil, Complexity.polyTimeComputable_pairFstD}
\difficulty{1}
For a bit string $s$, let $\pi_1$ be the total first projection from the self-delimiting
pair parser. The function returning the one-bit list whose bit is true exactly when
$\pi_1(s)$ is empty is polynomial-time computable.
\end{theorem}

\begin{theorem}[Polynomial time for the nested empty-head test]
\lean{Complexity.BlockLoop.polyTimeComputable_countdownNil}
\label{Complexity.BlockLoop.polyTimeComputable_countdownNil}
\leanok
\uses{Complexity.BlockLoop.polyTimeComputable_headNil, Complexity.PolyTimeComputable, Complexity.PolyTimeComputable.comp, Complexity.isNilB, Complexity.pairFstD, Complexity.pairSndD, Complexity.polyTimeComputable_pairSndD}
\difficulty{1}
For a bit string $s$, let $\pi_2$ be the total second projection from the
self-delimiting pair parser and $\pi_1$ its total first projection. The function
returning the one-bit list whose bit is true exactly when $\pi_1(\pi_2(s))$ is empty is
polynomial-time computable.
\end{theorem}

\begin{theorem}[Polynomial-time closure under a block-loop step]
\lean{Complexity.BlockLoop.polyTimeComputable_step}
\label{Complexity.BlockLoop.polyTimeComputable_step}
\leanok
\uses{Complexity.BlockLoop.polyTimeComputable_countdownNil, Complexity.BlockLoop.polyTimeComputable_headNil, Complexity.BlockLoop.step, Complexity.PolyTimeComputable, Complexity.PolyTimeComputable.comp, Complexity.PolyTimeComputable.pairEncode, Complexity.blockDone, Complexity.isNilB, Complexity.pairFstD, Complexity.pairSndD, Complexity.polyTimeComputable_const, Complexity.polyTimeComputable_ite, Complexity.polyTimeComputable_or, Complexity.polyTimeComputable_pairFstD, Complexity.polyTimeComputable_pairSndD, Complexity.polyTimeComputable_tail}
\difficulty{4}
Let $g$ be a function on finite bit strings, and let
$\mathrm{step}(\mathrm{upd},\mathrm{nxt})$ be the block-loop transformation defined by
the paired-string encoding. If $g=\mathrm{step}(\mathrm{upd},\mathrm{nxt})$, and both
functions $s\mapsto \mathrm{upd}(\pi_1(s),\pi_2(\pi_2(s)))$ and $s\mapsto
\mathrm{nxt}(\pi_1(s),\pi_2(\pi_2(s)))$ are computable in polynomial time, then $g$ is
computable in polynomial time.
\end{theorem}

\begin{theorem}[Polynomial time for block output]
\lean{Complexity.BlockLoop.polyTimeComputable_emit}
\label{Complexity.BlockLoop.polyTimeComputable_emit}
\leanok
\uses{Complexity.BlockLoop.emit, Complexity.BlockLoop.polyTimeComputable_countdownNil, Complexity.BlockLoop.polyTimeComputable_headNil, Complexity.PolyTimeComputable, Complexity.polyTimeComputable_const, Complexity.polyTimeComputable_ite}
\difficulty{3}
Let $\mathrm{out}$ be a polynomial-time computable function on bit strings. Define an
output function on bit strings $s$ as follows: return the empty string if the first
projection of $s$ is empty; otherwise return $\mathrm{out}(s)$ if the first projection
of the second projection of $s$ is empty; in all other cases return the empty string.
This function is polynomial-time computable.
\end{theorem}

\begin{theorem}[Polynomial time for countdown initialization]
\lean{Complexity.BlockLoop.polyTimeComputable_countdownInit}
\label{Complexity.BlockLoop.polyTimeComputable_countdownInit}
\leanok
\uses{Complexity.PolyTimeComputable, Complexity.PolyTimeComputable.comp, Complexity.PolyTimeComputable.pairEncode, Complexity.polyTimeComputable_const, Complexity.polyTimeComputable_polyUnary, Turing.pairEncode}
\difficulty{1}
Let $h_0$ be a bit string, let $x,r$ be polynomial-time computable functions on bit
strings, and let $a',k'\in\mathbb{N}$. Then the function sending $z$ to $\langle
h_0,\langle 1^{a'(|x(z)|+1)^{k'}},r(z)\rangle\rangle$ is polynomial-time computable.
\end{theorem}

\begin{theorem}[The completed block is fixed]
\lean{Complexity.BlockLoop.step_done}
\label{Complexity.BlockLoop.step_done}
\leanok
\uses{Complexity.BlockLoop.step, Complexity.blockDone, Complexity.isNilB}
\difficulty{1}
Let $g$ be a function on bit strings, and let $\mathrm{step}(\mathrm{upd},\mathrm{nxt})$
be the block-loop transformation. If $g=\mathrm{step}(\mathrm{upd},\mathrm{nxt})$, then
applying $g$ to the completed-block encoding returns that same completed-block encoding.
\end{theorem}

\begin{theorem}[One step bound for encoded states]
\lean{Complexity.BlockLoop.length_step_le}
\label{Complexity.BlockLoop.length_step_le}
\leanok
\uses{Complexity.BlockLoop.step, Complexity.blockDone, Complexity.eq_pairEncode_of_pairFstD_ne, Complexity.isNilB, Complexity.pairFstD, Complexity.pairFstD_nil, Complexity.pairFstD_pairEncode, Complexity.pairSndD, Complexity.pairSndD_pairEncode, Turing.length_pairEncode}
\difficulty{5}
Let $g$ be the step map determined by update and next-state maps, and let $c$ be a
natural number. Suppose that for every pair of finite bit strings $h,r$ with $h$
nonempty, twice the length of the updated string plus the length of the next-state
string is at most twice $|h|+|r|+2+c$. Then for every state string $s$, the length of
$g(s)$ plus $c$ times the length of the first projection of the second projection of
$g(s)$ is at most the maximum of $2$ and the corresponding quantity for $s$.
\end{theorem}

\begin{theorem}[Iterate length bound from slack]
\lean{Complexity.BlockLoop.length_iterate_le_of_slack}
\label{Complexity.BlockLoop.length_iterate_le_of_slack}
\leanok
\uses{Complexity.length_pairFstD_le, Complexity.length_pairSndD_le, Complexity.pairFstD, Complexity.pairSndD}
\difficulty{4}
Let $g$ be a map on finite bit strings and let $c,b$ be natural numbers with $2\le b$
and $c+1\le b$. Suppose that for every bit string $s$, the length of $g(s)$ plus $c$
times the length of the second component's first projection in $g(s)$ is at most the
maximum of $2$ and the corresponding quantity for $s$. Then for every bit string $w$ and
natural number $i$, the length of the $i$th iterate of $g$ from $w$ is at most
$b(|w|+1)$.
\end{theorem}

\begin{theorem}[Taking a slice after an encoded drop]
\lean{Complexity.BlockLoop.sliceTakeAt_pairEncode_drop}
\label{Complexity.BlockLoop.sliceTakeAt_pairEncode_drop}
\leanok
\uses{Complexity.pairFstD_pairEncode, Complexity.pairSndD_pairEncode, Complexity.sliceTakeAt, Turing.pairEncode}
\difficulty{1}
For natural numbers $a,k,i$ and bit strings $x,u$, let $L=a(|x|+1)^k$. Applying the
slice-take operation to the encoding of $x$ paired with $u$ after its first $iL$ bits
have been removed gives the encoding of $x$ paired with the first $L$ bits of that
remaining suffix (or the entire suffix if it has fewer than $L$ bits).
\end{theorem}

\begin{theorem}[Dropping a successive encoded slice]
\lean{Complexity.BlockLoop.sliceDropAt_pairEncode_drop}
\label{Complexity.BlockLoop.sliceDropAt_pairEncode_drop}
\leanok
\uses{Complexity.pairFstD_pairEncode, Complexity.pairSndD_pairEncode, Complexity.sliceDropAt, Turing.pairEncode}
\difficulty{3}
For natural numbers $a,k,i$ and bit strings $x,u$, let $L=a(|x|+1)^k$. Applying the
slice-drop operation to the encoding of $x$ paired with $u$ after its first $iL$ bits
have been removed gives the encoding of $x$ paired with $u$ after its first $(i+1)L$
bits have been removed.
\end{theorem}

\begin{theorem}[Orbit of a block countdown]
\lean{Complexity.BlockLoop.orbit}
\label{Complexity.BlockLoop.orbit}
\leanok
\uses{Complexity.BlockLoop.step, Complexity.isNilB, Turing.pairEncode}
\difficulty{5}
Let $P,Q:\mathbb{N}\to\{0,1\}^*$, let $K\in\mathbb{N}$, and let $s_0$ be a bit string.
Let $g$ be the block-loop step determined by update and next functions
$\mathrm{upd},\mathrm{nxt}:\{0,1\}^*\times\{0,1\}^*\to\{0,1\}^*$, and write $\langle
u,v\rangle$ for the self-delimiting pair encoding. Suppose $s_0=\langle P(0),\langle
1^K,Q(0)\rangle\rangle$, every $P(i)$ is nonempty for $i\le K$, and for every $i<K$,
$\mathrm{upd}(P(i),Q(i))=P(i+1)$ and $\mathrm{nxt}(P(i),Q(i))=Q(i+1)$. Then for each
$i\le K$, the $i$-th iterate of $g$ at $s_0$ is $\langle P(i),\langle
1^{K-i},Q(i)\rangle\rangle$.
\end{theorem}

\begin{theorem}[Orbit reaches the completed state]
\lean{Complexity.BlockLoop.iterate_done}
\label{Complexity.BlockLoop.iterate_done}
\leanok
\uses{Complexity.BlockLoop.step, Complexity.BlockLoop.step_done, Complexity.blockDone, Complexity.isNilB, Turing.pairEncode}
\difficulty{5}
Let $g$ be a map on finite bit strings that is the step map determined by update and
next-state maps. If $p$ is nonempty and the $K$th iterate of $g$ from $s_0$ is the pair
encoding of $p$ and the pair encoding of the empty string with $q$, then every iterate
with index $i>K$ equals the completed-block string.
\end{theorem}

\begin{theorem}[Strict polynomial budget bound]
\lean{Complexity.BlockLoop.lt_polyBudget_of_le}
\label{Complexity.BlockLoop.lt_polyBudget_of_le}
\leanok
\difficulty{2}
For natural numbers $a'$, $k'$, $n$, and $m$ with $n\le m$, one has
$a'(n+1)^{k'}<a'(m+1)^{k'}+1$.
\end{theorem}

\begin{theorem}[Orbit emission over a finite range]
\lean{Complexity.BlockLoop.flatMap_emit_orbit}
\label{Complexity.BlockLoop.flatMap_emit_orbit}
\leanok
\uses{Complexity.BlockLoop.emit, Complexity.blockDone, Complexity.flatMap_range_eq_single, Complexity.isNilB, Turing.pairEncode}
\difficulty{5}
Let $g$ and $\mathrm{out}$ be maps on finite bit strings, let
$P,Q:\mathbb{N}\to\{0,1\}^*$, and let $K,N$ be natural numbers with $K<N$. Suppose that
for every $i\le K$, the $i$th iterate of $g$ from $s_0$ is the pair encoding of $P(i)$
and the pair encoding of the string of $K-i$ ones with $Q(i)$; suppose also that every
iterate with index greater than $K$ equals the completed-block string, and that $P(i)$
is nonempty for all $i\le K$. Then concatenating the emissions at iterates indexed from
$0$ through $N-1$ gives precisely $\mathrm{out}$ applied to the pair encoding of $P(K)$
and the pair encoding of the empty string with $Q(K)$.
\end{theorem}

\begin{theorem}[Output of a finite block loop]
\lean{Complexity.BlockLoop.loop_output}
\label{Complexity.BlockLoop.loop_output}
\leanok
\uses{Complexity.BlockLoop.emit, Complexity.BlockLoop.flatMap_emit_orbit, Complexity.BlockLoop.iterate_done, Complexity.BlockLoop.orbit, Complexity.BlockLoop.step, Turing.pairEncode}
\difficulty{4}
Let $g$ and $\mathrm{out}$ be maps on finite bit strings, let $\mathrm{upd}$ and
$\mathrm{nxt}$ be maps on pairs of finite bit strings, let $P,Q:\mathbb{N}\to\{0,1\}^*$,
and let $K,N$ be natural numbers and $s_0$ a finite bit string. Suppose $s_0=\langle
P(0),\langle 1^K,Q(0)\rangle\rangle$, each $P(i)$ for $i\le K$ is nonempty, and for
every $i<K$, $\mathrm{upd}(P(i),Q(i))=P(i+1)$ and $\mathrm{nxt}(P(i),Q(i))=Q(i+1)$. If
$g$ is the step map determined by these update and next-state maps and $K<N$, then
concatenating the emissions at iterates indexed from $0$ through $N-1$ gives
$\mathrm{out}(\langle P(K),\langle [],Q(K)\rangle\rangle)$.
\end{theorem}

\begin{theorem}[Polynomial time for a bounded loop]
\lean{Complexity.BlockLoop.polyTimeComputable_of_loop}
\label{Complexity.BlockLoop.polyTimeComputable_of_loop}
\leanok
\uses{Complexity.PolyTimeComputable, Complexity.PolyTimeComputable.comp, Complexity.polyTimeComputable_emitIter}
\difficulty{3}
Let $g,e,\mathrm{init},f$ be functions on bit strings, and suppose each of
$g,e,\mathrm{init}$ is polynomial-time computable. Fix $a',k',b\in\mathbb{N}$, and
suppose that for every bit string $w$ and natural number $i$, the length of the $i$-th
iterate of $g$ at $w$ is at most $b(|w|+1)$. If, for every bit string $z$, concatenating
$e(g^{[i]}(\mathrm{init}(z)))$ over $0\le i<a'(|\mathrm{init}(z)|+1)^{k'}+1$ gives
$f(z)$, then $f$ is polynomial-time computable.
\end{theorem}

\begin{theorem}[Polynomial time for countdown output]
\lean{Complexity.BlockLoop.polyTimeComputable_countdown}
\label{Complexity.BlockLoop.polyTimeComputable_countdown}
\leanok
\uses{Complexity.BlockLoop.emit, Complexity.BlockLoop.length_iterate_le_of_slack, Complexity.BlockLoop.length_step_le, Complexity.BlockLoop.polyTimeComputable_countdownInit, Complexity.BlockLoop.polyTimeComputable_emit, Complexity.BlockLoop.polyTimeComputable_of_loop, Complexity.BlockLoop.polyTimeComputable_step, Complexity.BlockLoop.step, Complexity.PolyTimeComputable, Complexity.pairFstD, Complexity.pairSndD, Turing.pairEncode}
\difficulty{3}
Let $g,e,x,r,f$ be functions on bit strings, with $g$ the block-loop step determined by
$\mathrm{upd},\mathrm{nxt}$ and $e$ the corresponding output function determined by
$\mathrm{out}$. Fix a bit string $h_0$ and natural numbers $a',k',c$. Suppose the
functions $s\mapsto\mathrm{upd}(\pi_1(s),\pi_2(\pi_2(s)))$,
$s\mapsto\mathrm{nxt}(\pi_1(s),\pi_2(\pi_2(s)))$, $\mathrm{out}$, $x$, and $r$ are
polynomial-time computable, where $\pi_1,\pi_2$ are the total projections from the
self-delimiting pair parser. Suppose also that for every nonempty bit string $h$ and
every bit string $q$, $2|\mathrm{upd}(h,q)|+|\mathrm{nxt}(h,q)|\le 2|h|+|q|+2+c$. If,
for every bit string $z$, concatenating the outputs $e(g^{[i]}(\langle h_0,\langle
1^{a'(|x(z)|+1)^{k'}},r(z)\rangle\rangle))$ over all integers $i$ with $0\le
i<a'(|x(z)|+1)^{k'}+1$ gives $f(z)$, then $f$ is polynomial-time computable.
\end{theorem}

\begin{definition}[Pairwise XOR of decoded lists]
\lean{Complexity.xorD}
\label{Complexity.xorD}
\leanok
\uses{Complexity.pairFstD, Complexity.pairSndD}
For a Boolean list $z$, take the two lists obtained by decoding it as a pair, and return
their positionwise exclusive OR, stopping when either list ends.
\end{definition}

\begin{definition}[Encoded pair step]
\lean{Complexity.xorPairStep}
\label{Complexity.xorPairStep}
\leanok
\uses{Complexity.pairFstD, Complexity.pairSndD, Turing.pairEncode}
For a Boolean list $s$, decode it as a pair of lists, remove the first element of each
decoded list, and encode the resulting pair.
\end{definition}

\begin{definition}[Encoded pair emission]
\lean{Complexity.xorPairEmit}
\label{Complexity.xorPairEmit}
\leanok
\uses{Complexity.isNilB, Complexity.pairFstD, Complexity.pairSndD}
For a Boolean list $s$, decode it as a pair of lists. If either list is empty, return
the empty list; otherwise return a one-element list containing the exclusive OR of their
first elements.
\end{definition}

\begin{theorem}[Polynomial-time paired-tail transformation]
\lean{Complexity.polyTimeComputable_xorPairStep}
\label{Complexity.polyTimeComputable_xorPairStep}
\leanok
\uses{Complexity.PolyTimeComputable, Complexity.PolyTimeComputable.comp, Complexity.PolyTimeComputable.pairEncode, Complexity.polyTimeComputable_pairFstD, Complexity.polyTimeComputable_pairSndD, Complexity.polyTimeComputable_tail, Complexity.xorPairStep}
\difficulty{3}
For a bit string $s$, parse its two components using the aligned pair parser, taking
both components to be empty if parsing fails. The function that encodes as a pair the
tails obtained by deleting the first bit of each component is polynomial-time
computable.
\end{theorem}

\begin{theorem}[Polynomial-time first-bit XOR output]
\lean{Complexity.polyTimeComputable_xorPairEmit}
\label{Complexity.polyTimeComputable_xorPairEmit}
\leanok
\uses{Complexity.PolyTimeComputable, Complexity.PolyTimeComputable.comp, Complexity.isNilB, Complexity.pairFstD, Complexity.pairSndD, Complexity.polyTimeComputable_const, Complexity.polyTimeComputable_headD, Complexity.polyTimeComputable_isNil, Complexity.polyTimeComputable_ite, Complexity.polyTimeComputable_or, Complexity.polyTimeComputable_pairFstD, Complexity.polyTimeComputable_pairSndD, Complexity.xorPairEmit}
\difficulty{4}
For a bit string $s$, parse its two components using the aligned pair parser, taking
both components to be empty if parsing fails. The function that returns the one-bit word
given by the exclusive OR of the first bits of the two components when both components
are nonempty, and returns the empty word otherwise, is polynomial-time computable.
\end{theorem}

\begin{theorem}[Length bound for a pairing step]
\lean{Complexity.length_xorPairStep_le}
\label{Complexity.length_xorPairStep_le}
\leanok
\uses{Complexity.length_pairEncode_drop_le, Complexity.xorPairStep}
\difficulty{1}
For every bit string $s$, encode the strings obtained by removing the first bit from
each of its total pair projections. The resulting bit string has length at most
$\max(|s|,2)$.
\end{theorem}

\begin{theorem}[Orbit of the pair step]
\lean{Complexity.xorPairStep_orbit}
\label{Complexity.xorPairStep_orbit}
\leanok
\uses{Complexity.pairFstD_pairEncode, Complexity.pairSndD_pairEncode, Complexity.xorPairStep, Turing.pairEncode}
\difficulty{4}
For bit strings $p,u$ and an integer $i\ge 0$, applying the pair-step map $i$ times to
the encoding of $(p,u)$ gives the encoding of the pair obtained by deleting the first
$i$ bits of each component.
\end{theorem}

\begin{theorem}[XOR of the leading bits]
\lean{Complexity.xorPairEmit_cons}
\label{Complexity.xorPairEmit_cons}
\leanok
\uses{Complexity.isNilB, Complexity.xorPairEmit, Turing.pairEncode}
\difficulty{4}
For bits $b,c$ and bit strings $p,u$, apply the pair-output rule to the encoding of $b$
followed by $p$ and $c$ followed by $u$. The result is the one-bit string whose bit is
the XOR of $b$ and $c$.
\end{theorem}

\begin{theorem}[Empty output for an empty component]
\lean{Complexity.xorPairEmit_nil}
\label{Complexity.xorPairEmit_nil}
\leanok
\uses{Complexity.isNilB, Complexity.xorPairEmit, Turing.pairEncode}
\difficulty{4}
For bit strings $p,u$, if either $p$ or $u$ is empty, the pair-output rule applied to
their encoding returns the empty bit string.
\end{theorem}

\begin{theorem}[Output of the pairwise XOR process]
\lean{Complexity.xorPair_output}
\label{Complexity.xorPair_output}
\leanok
\uses{Complexity.xorPairEmit, Complexity.xorPairEmit_cons, Complexity.xorPairEmit_nil, Turing.pairEncode}
\difficulty{5}
For a natural number $N$ and bit strings $p,u$, if $\min(|p|,|u|)\le N$, then
concatenating the outputs of the pair-output rule on the encodings of the suffixes
obtained by deleting $i$ bits from each string, for $0\le i<N$, gives the bitwise XOR of
$p$ and $u$ up to the length of the shorter string.
\end{theorem}

\begin{theorem}[Polynomial-time decoded bitwise XOR]
\lean{Complexity.polyTimeComputable_xorD}
\label{Complexity.polyTimeComputable_xorD}
\leanok
\uses{Complexity.BlockLoop.length_iterate_le_of_slack, Complexity.BlockLoop.polyTimeComputable_of_loop, Complexity.PolyTimeComputable, Complexity.PolyTimeComputable.pairEncode, Complexity.length_pairFstD_le, Complexity.length_xorPairStep_le, Complexity.pairFstD, Complexity.pairSndD, Complexity.polyTimeComputable_pairFstD, Complexity.polyTimeComputable_pairSndD, Complexity.polyTimeComputable_xorPairEmit, Complexity.polyTimeComputable_xorPairStep, Complexity.xorD, Complexity.xorPairStep_orbit, Complexity.xorPair_output, Turing.length_pairEncode, Turing.pairEncode}
\difficulty{4}
For a bit string $z$, parse its two components using the aligned pair parser, taking
both components to be empty if parsing fails. The function that returns the bitwise
exclusive OR of the two components, up to the length of the shorter component, is
polynomial-time computable.
\end{theorem}

\begin{theorem}[XOR of an encoded pair]
\lean{Complexity.xorD_pairEncode}
\label{Complexity.xorD_pairEncode}
\leanok
\uses{Complexity.xorD, Turing.pairEncode}
\difficulty{1}
For any finite Boolean strings $a$ and $b$, applying the total XOR operation to their
self-delimiting pair encoding yields the bitwise XOR of $a$ and $b$, paired position by
position up to the length of the shorter string.
\end{theorem}
